\documentclass[11pt]{article}

% ---------------------------------------------------------------------------
%  Symmetric directional derivatives of probing a Tucker tensor train.
%
%  Compile with pdflatex (needs: amsmath, amssymb, amsthm, algorithm,
%  algpseudocode, booktabs, hyperref, geometry -- all standard).
%
%  This note derives, from a single Leibniz rule for Taylor jets, the efficient
%  sweeping algorithm that the backend module t3toolbox/backend/probe_derivatives.py
%  implements. The symmetric-derivative material is NOT part of the published T4S
%  paper (it was cut long ago); an older, unvetted derivation existed and is
%  superseded here. Everything below is verified against a dense oracle in
%  tests/test_probe_derivatives.py.
% ---------------------------------------------------------------------------

\usepackage[margin=1in]{geometry}
\usepackage{amsmath,amssymb,amsthm}
\usepackage{algorithm}
\usepackage{algpseudocode}
\usepackage{booktabs}
\usepackage[colorlinks=true,linkcolor=blue,citecolor=blue,urlcolor=blue]{hyperref}

\newcommand{\R}{\mathbb{R}}
\newcommand{\jet}[1]{[\![ #1 ]\!]}   % order-K jet (double square brackets, no stmaryrd dependency)
\newcommand{\dd}[1]{\,\mathrm{d}#1}
\DeclareMathOperator{\diag}{diag}

\theoremstyle{definition}
\newtheorem{definition}{Definition}
\newtheorem{remark}{Remark}
\theoremstyle{plain}
\newtheorem{lemma}{Lemma}
\newtheorem{theorem}{Theorem}
\newtheorem{corollary}{Corollary}

\title{\bfseries Symmetric directional derivatives of probing\\ a Tucker tensor train}
\author{T3Toolbox notes --- N.\ Alger, B.\ Christierson \&\ collaborators}
\date{\today}

\begin{document}
\maketitle

\begin{abstract}
Probing a Tucker tensor train (T3) evaluates its actions: contracting the tensor with a
collection of vectors in all modes but one. We derive the \emph{symmetric} directional
derivatives of this map --- the derivatives obtained by perturbing every input vector in the
\emph{same} direction $\mathbf P$, repeated $k$ times --- and show they are computed by a single
probing-style sweep that carries an extra Taylor-jet axis. The whole construction follows from one
Leibniz rule for jets: because each input vector enters affinely, its jet truncates at order one,
and every quantity built from it is assembled by a binomial convolution of jets. This note records
the math behind \texttt{t3toolbox/backend/probe\_derivatives.py}. It treats the Euclidean (plain-T3)
forward map; adjoint and Riemannian extensions are noted but not derived. This material is not part
of the published T4S paper.
\end{abstract}

\section{Setup: probing and its actions}
\label{sec:setup}

We recall just enough of the probing construction to fix notation; see the T4S paper for the full
development.

A \emph{Tucker tensor train} is a pair of core families $\mathbb T = \big((\mathbf U_i)_{i=1}^d,\,
(\mathbf G_i)_{i=1}^d\big)$, where each \emph{Tucker factor} $\mathbf U_i \in \R^{n_i \times N_i}$
has full row rank ($n_i \le N_i$) and each \emph{tensor-train core} $\mathbf G_i \in \R^{r_{i-1}
\times n_i \times r_i}$, with boundary ranks $r_0 = r_d = 1$. It represents a tensor $\mathbf T \in
\R^{N_1 \times \cdots \times N_d}$. For a vector $\mathbf v \in \R^{n_i}$ we write $\mathbf
G_i(\mathbf v) \in \R^{r_{i-1} \times r_i}$ for the matrix obtained by contracting the middle index,
\[
  \mathbf G_i(\mathbf v)[a,b] \;=\; \sum_{n} \mathbf G_i[a,n,b]\,\mathbf v[n].
\]

\paragraph{Actions (probing).}
Given input vectors $\mathbf W = (\mathbf w_1, \dots, \mathbf w_d)$ with $\mathbf w_i \in \R^{N_i}$,
the $i$th \emph{action} $\mathbf z_i$ is the result of contracting $\mathbf T$ with $\mathbf w_j$ in
every mode $j \ne i$, leaving mode $i$ free:
\[
  \mathbf z_i \;=\; T(\mathbf w_1, \dots, \mathbf w_{i-1}, \,\cdot\,, \mathbf w_{i+1}, \dots, \mathbf w_d)
  \;\in\; \R^{N_i}.
\]
Probing computes all $d$ actions in two sweeps over the train. Introduce the \emph{edge variables}
\begin{align}
  \boldsymbol\xi_i      &= \mathbf U_i \mathbf w_i \in \R^{n_i},
  & &\text{(projected probe)} \label{eq:xi}\\
  \boldsymbol\mu_0^\top  &= 1, &
  \boldsymbol\mu_i^\top &= \boldsymbol\mu_{i-1}^\top\, \mathbf G_i(\boldsymbol\xi_i) \in \R^{r_i},
  & &\text{(left pushthrough)} \label{eq:mu}\\
  \boldsymbol\nu_d      &= 1, &
  \boldsymbol\nu_i      &= \mathbf G_{i+1}(\boldsymbol\xi_{i+1})\, \boldsymbol\nu_{i+1} \in \R^{r_i},
  & &\text{(right pushthrough)} \label{eq:nu}\\
  \boldsymbol\eta_i     &= \boldsymbol\mu_{i-1}^\top\, \mathbf G_i\, \boldsymbol\nu_i \in \R^{n_i},
  & &\text{(combine, mode $i$ free)} \label{eq:eta}\\
  \mathbf z_i           &= \mathbf U_i^\top \boldsymbol\eta_i \in \R^{N_i}.
  & &\text{(lift)} \label{eq:z}
\end{align}
Here $\boldsymbol\mu_{i-1}^\top \mathbf G_i \boldsymbol\nu_i$ denotes the vector obtained by
contracting $\boldsymbol\mu_{i-1}$ and $\boldsymbol\nu_i$ into the left and right bond indices of
$\mathbf G_i$, leaving the middle index free. The single Tucker factor $\mathbf U_i$ is used twice:
to \emph{project} the probe \eqref{eq:xi} and to \emph{lift} the result \eqref{eq:z}.

\paragraph{Goal.}
Fix a direction $\mathbf P = (\mathbf p_1, \dots, \mathbf p_d)$, $\mathbf p_i \in \R^{N_i}$. We seek
the \emph{symmetric directional derivatives} of each action with respect to the inputs, taken $k$
times in the single repeated direction $\mathbf P$:
\begin{equation}
  \label{eq:goal}
  \mathbf z_i^{(k)} \;:=\; \frac{\dd{^k}}{\dd s^k}\, \mathbf z_i(\mathbf W + s\mathbf P)\Big|_{s=0}
  \;=\; D^k \mathbf z_i(\mathbf W)\,\mathbf P^{k},
  \qquad k = 0, 1, \dots, K.
\end{equation}
Order $k=0$ is the ordinary probe. The point is that the \emph{same} direction is fed into all $k$
slots; this is what makes the whole family $k=0,\dots,K$ computable in a single sweep.

\section{Jets and the Leibniz rule}
\label{sec:jets}

The entire derivation rests on viewing every quantity as a function of the scalar $s$ along the ray
$\mathbf W + s\mathbf P$ and tracking its Taylor data.

\begin{definition}[Jet]
\label{def:jet}
Let $q$ be a quantity that depends smoothly on the inputs $\mathbf W$. Its \emph{order-$K$ jet}
(along $\mathbf P$, at $\mathbf W$) is the tuple of symmetric derivatives
\[
  \jet{q} \;=\; \big(q^{(0)}, q^{(1)}, \dots, q^{(K)}\big),
  \qquad
  q^{(k)} := \frac{\dd{^k}}{\dd s^k}\, q(\mathbf W + s\mathbf P)\Big|_{s=0}.
\]
So $q^{(0)} = q(\mathbf W)$ is the base value and $q^{(k)}$ is the quantity \eqref{eq:goal} sought
for $q = \mathbf z_i$.
\end{definition}

\begin{lemma}[Affine inputs have order-one jets]
\label{lem:affine}
Each probe vector is affine in $s$, $\mathbf w_i(s) = \mathbf w_i + s\,\mathbf p_i$, so
\[
  \jet{\mathbf w_i} = (\mathbf w_i,\ \mathbf p_i,\ 0, \dots, 0).
\]
Consequently the projected probe \eqref{eq:xi}, being a fixed linear image, also has an order-one
jet:
\[
  \jet{\boldsymbol\xi_i} = (\boldsymbol\xi_i,\ \delta\boldsymbol\xi_i,\ 0, \dots, 0),
  \qquad \delta\boldsymbol\xi_i := \mathbf U_i \mathbf p_i.
\]
\end{lemma}

\begin{proof}
$\mathbf w_i(s)$ is affine, so its second and higher $s$-derivatives vanish. The projection
$\boldsymbol\xi_i(s) = \mathbf U_i \mathbf w_i(s) = \mathbf U_i \mathbf w_i + s\,\mathbf U_i \mathbf
p_i$ is again affine.
\end{proof}

\begin{lemma}[Leibniz rule for jets]
\label{lem:leibniz}
Let $B : \R^{A} \times \R^{B} \to \R^{C}$ be a fixed bilinear map, and let $a(s), b(s)$ depend
smoothly on $s$. Then the jet of $s \mapsto B(a(s), b(s))$ is the binomial convolution of the jets of
$a$ and $b$:
\begin{equation}
  \label{eq:leibniz}
  \big[B(a,b)\big]^{(k)} \;=\; \sum_{j=0}^{k} \binom{k}{j}\, B\!\big(a^{(j)},\, b^{(k-j)}\big),
  \qquad k = 0, 1, \dots, K.
\end{equation}
\end{lemma}

\begin{proof}
By bilinearity, $\tfrac{\dd}{\dd s} B(a,b) = B(a',b) + B(a,b')$. This is the ordinary product rule
for the bilinear pairing $B$, so the general Leibniz rule for the $k$th derivative of a product gives
\eqref{eq:leibniz}; the binomial coefficients count how the $k$ differentiations are distributed
between the two arguments. Evaluating at $s=0$ replaces each derivative by the corresponding jet
component.
\end{proof}

Lemma~\ref{lem:leibniz} is the only engine we need. Every edge variable in
\eqref{eq:xi}--\eqref{eq:z} is built from the inputs by a chain of bilinear pairings (and one final
linear map), so its jet is obtained by applying \eqref{eq:leibniz} at each pairing. The two regimes
that occur differ only in how many jet components the second argument has: \emph{the input jet has
just two} (Lemma~\ref{lem:affine}), while \emph{a pushthrough jet has all $K{+}1$}.

\begin{remark}[Jets as a truncated algebra]
\label{rem:algebra}
Equivalently, identify $\jet{q}$ with the truncated Taylor polynomial $\sum_{k} \tfrac{1}{k!}
q^{(k)} s^{k} \pmod{s^{K+1}}$. Then \eqref{eq:leibniz} is exactly multiplication in the algebra
$\R[s]/(s^{K+1})$ (the Cauchy product, with the $\tfrac1{k!}$ normalisation absorbed into the
binomial weights). We keep the unnormalised derivatives $q^{(k)}$ rather than Taylor coefficients
because the binomial form maps directly onto integer-exact arithmetic in the implementation.
\end{remark}

\section{Symmetric derivatives of the actions}
\label{sec:derivs}

We now apply Lemma~\ref{lem:leibniz} to each edge-variable recursion. Throughout, the cores $\mathbf
U_i, \mathbf G_i$ are constants (independent of $s$), so the only $s$-dependence is carried by the
edge variables.

\begin{theorem}[Pushthrough jets]
\label{thm:push}
The left and right pushthrough jets satisfy, for $k = 1, \dots, K$,
\begin{align}
  \big(\boldsymbol\mu_i^{(k)}\big)^\top
    &= \big(\boldsymbol\mu_{i-1}^{(k)}\big)^\top \mathbf G_i(\boldsymbol\xi_i)
     \;+\; k\,\big(\boldsymbol\mu_{i-1}^{(k-1)}\big)^\top \mathbf G_i(\delta\boldsymbol\xi_i),
    \label{eq:dmu}\\[2pt]
  \boldsymbol\nu_i^{(k)}
    &= \mathbf G_{i+1}(\boldsymbol\xi_{i+1})\, \boldsymbol\nu_{i+1}^{(k)}
     \;+\; k\,\mathbf G_{i+1}(\delta\boldsymbol\xi_{i+1})\, \boldsymbol\nu_{i+1}^{(k-1)},
    \label{eq:dnu}
\end{align}
with base values $\boldsymbol\mu_i^{(0)} = \boldsymbol\mu_i$, $\boldsymbol\nu_i^{(0)} =
\boldsymbol\nu_i$, and boundary jets $\boldsymbol\mu_0^{(0)} = 1$, $\boldsymbol\mu_0^{(k)} = 0$
($k>0$), $\boldsymbol\nu_d^{(0)} = 1$, $\boldsymbol\nu_d^{(k)} = 0$ ($k>0$).
\end{theorem}

\begin{proof}
The recursion \eqref{eq:mu} reads $\boldsymbol\mu_i^\top = B(\boldsymbol\mu_{i-1},
\boldsymbol\xi_i)$ with the bilinear map $B(\mathbf m, \mathbf v) = \mathbf m^\top \mathbf G_i(\mathbf
v)$ (linear in $\mathbf m$, and linear in $\mathbf v$ since $\mathbf G_i(\cdot)$ is). Apply
Lemma~\ref{lem:leibniz}:
\[
  \big(\boldsymbol\mu_i^{(k)}\big)^\top
  = \sum_{j=0}^{k} \binom{k}{j}\, \big(\boldsymbol\mu_{i-1}^{(j)}\big)^\top
      \mathbf G_i\!\big(\boldsymbol\xi_i^{(k-j)}\big).
\]
By Lemma~\ref{lem:affine}, $\boldsymbol\xi_i^{(k-j)}$ is nonzero only for $k-j \in \{0,1\}$, i.e.\
$j \in \{k, k-1\}$, leaving the two terms of \eqref{eq:dmu} with weights $\binom{k}{k}=1$ and
$\binom{k}{k-1}=k$. The boundary jet $\jet{\boldsymbol\mu_0}$ is constant ($\boldsymbol\mu_0 \equiv
1$), so all its derivatives above order $0$ vanish. Equation \eqref{eq:dnu} is identical with the
train reversed.
\end{proof}

Equation~\eqref{eq:dmu} is a \emph{bidiagonal} recursion in the order index $k$: the order is
preserved when the core is fed the base projection $\boldsymbol\xi_i$ and raised by one (with weight
$k$) when it is fed the direction $\delta\boldsymbol\xi_i$. The bidiagonal (only two-term) shape is
exactly the statement that the input jet stops at order one.

\begin{theorem}[Action jets]
\label{thm:action}
For $i = 1, \dots, d$ and $k = 0, 1, \dots, K$,
\begin{align}
  \boldsymbol\eta_i^{(k)}
    &= \sum_{j=0}^{k} \binom{k}{j}\,
       \big(\boldsymbol\mu_{i-1}^{(j)}\big)^\top \mathbf G_i\, \boldsymbol\nu_i^{(k-j)},
    \label{eq:deta}\\[2pt]
  \mathbf z_i^{(k)} &= \mathbf U_i^\top \boldsymbol\eta_i^{(k)}.
    \label{eq:dz}
\end{align}
\end{theorem}

\begin{proof}
The combine \eqref{eq:eta} is $\boldsymbol\eta_i = B(\boldsymbol\mu_{i-1}, \boldsymbol\nu_i)$ with the
bilinear map $B(\mathbf m, \mathbf n) = \mathbf m^\top \mathbf G_i \mathbf n$ (mode $i$ free), so
\eqref{eq:deta} is Lemma~\ref{lem:leibniz} verbatim. Unlike the pushthrough, \emph{both} arguments
carry full jets, so the sum runs over all $j = 0, \dots, k$. Finally \eqref{eq:z} is the fixed
linear map $\boldsymbol\eta_i \mapsto \mathbf U_i^\top \boldsymbol\eta_i$, whose jet acts
componentwise: $\mathbf z_i^{(k)} = \mathbf U_i^\top \boldsymbol\eta_i^{(k)}$.
\end{proof}

\begin{remark}[The full sum, $j = 0$ to $k$]
\label{rem:fullsum}
The action sum \eqref{eq:deta} runs over the \emph{full} range $j = 0, \dots, k$; both endpoints
($\boldsymbol\mu_{i-1}^{(0)} = \boldsymbol\mu_{i-1}$ paired with $\boldsymbol\nu_i^{(k)}$, and
$\boldsymbol\mu_{i-1}^{(k)}$ paired with $\boldsymbol\nu_i^{(0)} = \boldsymbol\nu_i$) are genuine
product-rule terms. An earlier unpublished draft wrote this sum starting at $j=1$; that is an error.
The forms above are verified against a dense oracle (Section~\ref{sec:impl}).
\end{remark}

\section{The sweeping algorithm}
\label{sec:algo}

Theorems~\ref{thm:push}--\ref{thm:action} give an algorithm that produces the entire derivative
family $\{\mathbf z_i^{(k)}\}_{i,k}$ in two sweeps, exactly mirroring ordinary probing but carrying
the order index $k$.

\begin{algorithm}[h]
\caption{Symmetric probe derivatives of a Tucker tensor train}
\label{alg:sweep}
\begin{algorithmic}[1]
\Require T3 $\big((\mathbf U_i),(\mathbf G_i)\big)$; inputs $\mathbf W$; direction $\mathbf P$; order $K$
\Ensure derivative jets $\jet{\mathbf z_i}$ for $i=1,\dots,d$
\For{$i = 1,\dots,d$} \Comment{project, forming the order-one input jets}
  \State $\boldsymbol\xi_i \gets \mathbf U_i \mathbf w_i$,\quad
         $\delta\boldsymbol\xi_i \gets \mathbf U_i \mathbf p_i$
\EndFor
\State $\jet{\boldsymbol\mu_0} \gets (1,0,\dots,0)$
\For{$i = 1,\dots,d$} \Comment{left sweep, Eq.~\eqref{eq:dmu}}
  \State $\boldsymbol\mu_i^{(k)} \gets \boldsymbol\mu_{i-1}^{(k)}\mathbf G_i(\boldsymbol\xi_i)
         + k\,\boldsymbol\mu_{i-1}^{(k-1)}\mathbf G_i(\delta\boldsymbol\xi_i)$ \textbf{ for } $k=0,\dots,K$
\EndFor
\State $\jet{\boldsymbol\nu_d} \gets (1,0,\dots,0)$
\For{$i = d,\dots,1$} \Comment{right sweep, Eq.~\eqref{eq:dnu}}
  \State $\boldsymbol\nu_{i-1}^{(k)} \gets \mathbf G_i(\boldsymbol\xi_i)\boldsymbol\nu_i^{(k)}
         + k\,\mathbf G_i(\delta\boldsymbol\xi_i)\boldsymbol\nu_i^{(k-1)}$ \textbf{ for } $k=0,\dots,K$
\EndFor
\For{$i = 1,\dots,d$} \Comment{combine and lift, Eqs.~\eqref{eq:deta}--\eqref{eq:dz}}
  \State $\boldsymbol\eta_i^{(k)} \gets \sum_{j=0}^{k}\binom{k}{j}
         \big(\boldsymbol\mu_{i-1}^{(j)}\big)^\top \mathbf G_i\,\boldsymbol\nu_i^{(k-j)}$ \textbf{ for } $k=0,\dots,K$
  \State $\mathbf z_i^{(k)} \gets \mathbf U_i^\top \boldsymbol\eta_i^{(k)}$ \textbf{ for } $k=0,\dots,K$
\EndFor
\end{algorithmic}
\end{algorithm}

\paragraph{Cost.}
Relative to a single probe, the pushthrough sweeps carry an order axis of length $K+1$, each step
combining only the orders $k$ and $k-1$ (the input jet is order one), so they cost a factor $O(K)$
more. The combine \eqref{eq:deta} is a length-$(K{+}1)$ binomial convolution at each free mode, an
$O(K^2)$ factor over a single combine. There are no other changes: the algorithm is the ordinary
probing pipeline with the order axis threaded through.

\section{The unifying view and the implementation}
\label{sec:impl}

\paragraph{One operation.}
Theorems~\ref{thm:push} and \ref{thm:action} are the \emph{same} statement --- Lemma~\ref{lem:leibniz}
--- applied to two different bilinear pairings. Both are driven by the binomial tensor
\begin{equation}
  \label{eq:trs}
  \texttt{trs}[t,r,s] \;=\; \binom{t}{r}\,[\,r + s = t\,],
\end{equation}
where $[\,\cdot\,]$ is the Iverson bracket. The pushthrough is the convolution of a pushthrough jet
(index $r$) with the input jet (index $s$, supported on $s \in \{0,1\}$), and the combine is the
convolution of the left and right pushthrough jets (indices $r$ and $s$, both full). In both cases
$r$ and $s$ are summed against $\texttt{trs}$ to produce the output order $t$. This is the product in
the truncated jet algebra of Remark~\ref{rem:algebra}.

\paragraph{Contractions.}
The implementation (\texttt{t3toolbox/backend/probe\_derivatives.py}) realises each step as a single
einsum, named by its indices (lowercase $t,r,s$ for the order axes, capital $W$ for the sample stack
and $C$ for the base stack, as in plain probing):
\begin{center}
\begin{tabular}{lll}
\toprule
step & contraction & probing analogue \\
\midrule
pushthrough \eqref{eq:dmu}/\eqref{eq:dnu}
  & \texttt{trs\_rWCa\_Caib\_sWCi\_to\_tWCb} & \texttt{WCa\_Caib\_WCi\_to\_WCb} \\
combine \eqref{eq:deta}
  & \texttt{trs\_rWCa\_Caib\_sWCb\_to\_tWCi} & \texttt{WCa\_Caib\_WCb\_to\_WCi} \\
lift \eqref{eq:dz}
  & \texttt{tWCi\_Cio\_to\_tWCo} & \texttt{WCi\_Cio\_to\_WCo} \\
\bottomrule
\end{tabular}
\end{center}
The first two share the tensor \eqref{eq:trs}; for the pushthrough its $s$ axis is truncated to
$\{0,1\}$. The order axis $t$ is outermost. A flat \emph{sample stack} lets one call evaluate many
$(\mathbf W, \mathbf P)$ pairs at once: each sample is a paired input/direction, so to sweep many
directions at one point one repeats that point across the stack.

\paragraph{Verification.}
Each action is multilinear in the off-mode inputs, so $\mathbf z_i(\mathbf W + s\mathbf P)$ is a
polynomial in $s$ and the $k$th derivative is exact:
\begin{equation}
  \label{eq:oracle}
  \mathbf z_i^{(k)} \;=\; k!\!\!\sum_{\substack{S \subseteq \{1,\dots,d\}\setminus\{i\}\\ |S| = k}}
  T\big(\dots\big),
\end{equation}
the inner contraction feeding $\mathbf p_j$ in the modes $j \in S$ and $\mathbf w_j$ elsewhere. The
sweeping algorithm matches this dense subset-expansion oracle to machine precision in
\texttt{tests/test\_probe\_derivatives.py}, across orders, ranks, and sample-stack shapes.

\section{Extensions (not derived here)}
\label{sec:future}

Two extensions follow the same jet calculus and are left for a future revision.

\emph{Adjoint and corewise derivatives.} The transpose of the (linear-in-the-tensor) derivative map,
and the gradient with respect to the cores treated as independent variables, are again binomial jet
convolutions --- of the adjoint edge variables against the same tensor \eqref{eq:trs}. These mirror
the three transpose flavours (ambient, corewise, tangent) already present for ordinary probing.

\emph{Riemannian derivatives.} Replacing the plain cores by an orthogonal frame and a tangent
variation gives the symmetric derivatives of \emph{tangent-vector} probing. The pushthrough and
combine acquire the variation terms of the tangent calculus, but the order axis threads through them
exactly as above; the jet Leibniz rule is unchanged.

\end{document}
